\section{Our Evaluation}\label{sec:evaluation}

\paragraph{Dataset.}
CVE descriptions can be very repetitive: it is common to observe several near-duplicate descriptions of the same vulnerability (e.g., found in different packages), sometimes weeks or months apart.
Thus, a random split of the CVE data can produce a training set which is representative of the test data distribution.
In this work, we instead take a temporal approach: train only on 2022 and 2023 CVEs and test on 2024 CVEs, to determine how well each system can generalize to future vulnerabilities.
All CVE data is taken from a particular snapshot of a National Vulnerability Database (NVD) feed,\footnote{Release \texttt{v2026.10.05-000016} of the Fraunhofer FKIE mirror of the NVD JSON feeds~\citep{fkie2026nvdfeeds}.} so as to avoid discrepancies between statistics reported in this document and results of released artifacts.
We evaluate on all \nTest{} 2024 CVEs meeting the following inclusion criteria (see Appendix~\ref{app:stats} for statistics and additional information):
\begin{itemize}
  \item CVE has a CVSS v3.x base metric vector and severity label;
  \item CVE has a CWE that is specific (excluding the placeholders \texttt{NVD-CWE-noinfo} and \texttt{NVD-CWE-Other});
  \item CVE's NVD-provided description is in English and at least 20 characters long.
\end{itemize}

\paragraph{Approach and metrics.}
As discussed in the introduction, CVSS base score and severity are simple calculations~\citep{first2019cvss31}, functions of the eight different components.
Thus, given a vulnerability description, we want to train systems that either (1) predict a score, or (2) predict the eight components and then calculate the score from that vector~\citep{elbaz2020fighting,shahid2021cvssbert}, and analyze performance in terms of severity classification (accuracy and macro-averaged F$_1$-score), CWE classification (accuracy), and CVSS score and vector.
For CVSS score, we report mean absolute error (MAE), and for CVSS vector, we report the percent of CVEs for which all eight components are correctly predicted (``Exact CVSS'' in Table~\ref{tab:summary}).
We implemented a CVSS v3.1 calculator and verified it agrees with the CVSS base score listed by the NVD for all \nAll{} CVEs in our dataset.
Thus, any score errors reported are a result of mistakes in component prediction.

\paragraph{Systems evaluated.}
For CWE classification, CVSS score, vector, and severity classification we evaluate three systems: a TF-IDF + logistic regression classifier, an instruction-tuned 3.8B-parameter LLM prompted zero-shot for classification, and that same LLM after a short (\nLoraTrainHours{} hour) LoRA fine-tuning to predict CWE and CVSS.
TF-IDF-based classifiers are strong baselines on this task: \citet{talibzade2025llms} find TF-IDF features perform best for CWE and severity prediction, and find them challenging to outperform on NLP tasks using vulnerability description text.
Further details about data preparation and training can be found in Appendix~\ref{app:data}.

We release our code\footnote{\url{https://github.com/jessiedong01/cvss-cwe-prediction}} as well as our system's prediction for each CVE.\footnote{In the \texttt{results/} directory of the repository, one JSON Lines file per system.}
Statistics reported here have been updated to be calculated from a particular snapshot of an NVD feed~\citep{fkie2026nvdfeeds} release (\texttt{v2026.10.05-000016}); please run our provided scripts for calculation on this data release to exactly reproduce the numbers below and in Appendix~\ref{app:results}.
